\section{Numerical experiments}
\label{sec:experiments}

Everything above is a proof, and proofs leave two things open. The first is whether
the code satisfies the identities: they hold in exact arithmetic, but the
implementation reaches them through a triangular factor held as packed columns
addressed by hand and through library calls writing in place into a view of $J$'s
own buffer, correct only while that view stays contiguous. A sign error in the
packed indexing would leave every proposition above true and the running solver
wrong, and wrong quietly: the invariants are maintained across a run, not
recomputed, and a solver drifting from $J\T G J = I$ still returns a vector that is
very nearly right. The second is the set of quantities the paper reasons about but
cannot predict: how often a guessed active set is certifiable, how many repairs it
takes, what a warm-started solve costs. All measurements are on the implementation
at the version recorded in Section~\ref{sec:reproducibility}, by scripts
distributed with this paper's sources.

\subsection{The guessed active set, and what the certificate catches}
\label{ssec:pdas}

Section~\ref{sec:pdas} makes two claims only measurement settles: that the block
exchange converges in a number of repairs independent of $n$, and that losing the
$P$-matrix property for general $C$ is a real exposure.
Figure~\ref{fig:pdas} addresses both over \qpPdasFamilies{} constraint families,
\qpPdasInstances{} instances each, at $n \in \{\qpPdasSizes\}$, instrumenting the
shipped routines instead of reimplementing them.

\begin{figure}[ht]
\centering
\includegraphics[width=0.86\textwidth]{graphs/quadprog_pdas.pdf}
\caption{Left: the fraction of guesses the certificate accepts. Right: set repairs
to convergence. Four families are certified without exception at every size and
settle in \qpPdasRepairsLo{} to \qpPdasRepairsHi{} repairs while the exact walk's
outer iterations grow linearly in $n$. The fifth is constructed so that a guessed
set can be linearly dependent, and it fails.}
\label{fig:pdas}
\end{figure}

On the four well-posed families (bounds, budget-plus-bounds, dense $C$, and $C$
with equalities) every guess is certified at every size, and the repair
count rises only from \qpPdasRepairsLo{} to \qpPdasRepairsHi{} across a
sixteen-fold range of $n$, never exceeding \qpPdasRepairsMax{}. The comparison
that matters is against the exact walk on the same instances: on
budget-plus-bounds it takes \qpPdasBudgetOuter{} outer iterations at
$n = \qpPdasRefN{}$, and its count grows linearly where the repair count is very
nearly flat. Section~\ref{sec:pdas} describes that trade; here it is measured.

\paragraph{The rank-deficient guess is real, and invisible by accident.} Over
\qpPdasInstances{} instances at each of five sizes, not one guess on those four
families produced a rank-deficient working-set system, from which a reader might
conclude the exposure is theoretical. It is not, merely invisible in random data,
where a dense $C$ almost never has a dependent subset the guess happens to name.
The fifth family makes it reachable in the most elementary way, by repeating
columns of $C$, so that naming both copies induces a singular system: there the
rank test fires on \qpPdasSingularDup\% of attempts and the certified fraction
collapses from \qpPdasCertSmall\% at $n = \qpPdasNsmall{}$ to \qpPdasCertBig\% at
$n = \qpPdasNbig{}$, against \qpPdasSingularBox\% on the others. Duplicated
constraints are not exotic, arising whenever a model is assembled programmatically
from overlapping rules, and they give the structural point of
Section~\ref{sec:pdas} its concrete form: for bounds the system is a principal
submatrix of $G^{-1}$ and cannot be singular, and nothing about general $C$
preserves that.

\paragraph{Where the failures are caught.} Of the candidates reaching the
certificate, \qpPdasAccepted{} were accepted with a worst KKT
residual~\eqref{eq:certificate} of $\qpPdasAcceptWorst$, and \qpPdasRejected{}
were rejected. We report that asymmetry instead of smoothing it: on these
families the certificate never had to reject a converged candidate, because the
failures were caught earlier and more cheaply by the Cholesky factorisation
refusing a dependent guess. The certificate remains the guarantee, but on this
evidence the rank test does the work in practice.

\subsection{The same program in worse coordinates}
\label{ssec:transform}

Proposition~\ref{prop:reparam} makes an experiment available that the rest of
this section cannot run. Substituting $x = Ty$ for a random invertible $T$ leaves
the active set, the multipliers and the sequence of working sets exactly where
they were, while $\kappa(G') \approx \kappa(G)\kappa(T)^2$ becomes whatever
$\kappa(T)$ is asked to make it. The transformed program therefore has a known
answer at an arbitrarily bad conditioning, obtained from a solve at a good one
rather than from a second implementation, which is what
Section~\ref{ssec:limits} otherwise has to do without. Two base problems, bounds
and dense $C$, at $n = \qpTfN{}$, with $\qpTfInstances{}$ instances at each of
eight values of $\kappa(T)$ up to $\qpTfKappaTHi$: $\qpTfSolves{}$ transformed
solves in all, reaching $\kappa(G') = \qpTfKappaHi$.

\begin{figure}[ht]
\centering
\includegraphics[width=0.86\textwidth]{graphs/quadprog_transform.pdf}
\caption{Left: the error in $x$ against the conditioning the transformation
produces, against the $\kappa(G')\varepsilon$ a backward-stable solve is
entitled to, with the band in which the certificate begins to reject shaded. The
active set is the reference's in every one of these solves, including those at
the right-hand end. Right: three policies on a family of
$\qpTfFamilyTerms{}$ programs sharing an ill-conditioned $G$.}
\label{fig:transform}
\end{figure}

\paragraph{The invariance survives the arithmetic; the coordinates do not.} The
active set agreed with the reference's on \qpTfSetAgree\% of the
$\qpTfSolves{}$ solves and the outer iteration count on \qpTfIterAgree\%, which
is the proposition reproduced --- and reproduced well past the point where the
answer itself has decayed. At the worst conditioning measured the error in $x$
is $\qpTfErrHi$ and in $\lambda$ is $\qpTfLamErrHi$, some two or three correct
digits, and the set of binding constraints is still exactly right. Combinatorial
output is the robust part of what this method returns, which for a long-only
portfolio is the statement that the holdings are identified long after their
sizes are not.
Figure~\ref{fig:transform} also confirms that the errors sit at or below
$\kappa(G')\varepsilon$ throughout: this is the transformation degrading the
problem, not the solver degrading with it.

\paragraph{The certificate's rejecting population, at last.} Section~\ref{ssec:pdas}
had nothing to report on the rejecting side, because on well-conditioned data the
rank test caught every failure first. Here the other case appears, and it
appears as a band rather than at a threshold: the rejected fraction across the
grid runs $\qpTfRates$\%, the first rejection arriving at
$\kappa(G') = \qpTfFirstBad$ while guesses are still being certified as far out
as $\qpTfLastGood$, a majority failing only from $\qpTfMajority$, and the two
bases disagreeing throughout the transition. The accepted candidates keep the
residual the certificate demands of them, $\qpTfAcceptWorst$ at worst. What is
being rejected is not a wrong active set, since the set is right throughout, but
a primal--dual pair the certificate cannot verify to tolerance, and it is right
to refuse: past this band the walk's own answer is no better, and nothing in
double precision recovers what the transformation has destroyed.
The remark of Section~\ref{sec:pdas} on the two tolerances should be read with
this in view. The separation it relies on is real, and it is a property of
well-conditioned data rather than of the certificate.

\paragraph{What a rejected guess costs.} At $n = \qpTfTimeN{}$ on the
well-conditioned end the fast path takes $\qpTfFastWell$\,ms against the exact
walk's $\qpTfWalkWell$\,ms, a margin of $\qpTfWellSpeedup\times$ on this
family and size, and of the same character as the ones
Section~\ref{ssec:compare} reports. Inside the rejecting band it takes
$\qpTfFastIll$\,ms against $\qpTfWalkIll$\,ms: $\qpTfIllPenalty$\% \emph{slower}
than not guessing, the guess being paid for and thrown away. So the ordering
Section~\ref{ssec:compare} reports has a conditioning boundary, and
Section~\ref{ssec:whennot} asks what to do about it. The two mistakes cost
$\qpTfFalseSkipCost$\,ms and $\qpTfMissedSkipCost$\,ms, and the free conditioning
estimate does not order them: $\qpTfEstGood$ at the worst accepted guess against
$\qpTfEstBad$ at the best rejected one. Remembering a rejection works where
predicting one does not. On a family of $\qpTfFamilyTerms{}$ linear terms over
one ill-conditioned $G$, where $\qpTfFamilyRejected{}$ of the
$\qpTfFamilyTerms{}$ guesses are rejected, guessing every time costs
$\qpTfFamilyAlways$\,ms, stopping after the first rejection costs
$\qpTfFamilyCached$\,ms, and never guessing costs $\qpTfFamilyNever$\,ms: the
rule recovers $\qpTfFamilySaving$\% of the waste and leaves a percent of the
floor on the table, in exchange for keeping the whole speedup on families where
guessing works.

\subsection{Against other solvers, and against the same method done differently}
\label{ssec:compare}

The experiments so far compare the solver with its own claims, which establishes
that it is implemented correctly and says nothing about whether the method is worth
using. Table~\ref{tab:compare} and Figure~\ref{fig:compare} answer that against
four alternatives. Two are outside the family: an operator-splitting method (OSQP)
and an interior-point method (Clarabel). Two are inside it, and are the
informative ones: the reference C implementation of this exact algorithm, and
DAQP~\cite{arnstrom2022daqp}, which walks the same dual active set but carries it
on recursive $LDL^\top$ updates of the working-set system rather than on the
$(J,R)$ pair of Section~\ref{sec:factorisation}.

Times alone would mislead, in the direction that flatters whichever solver was
configured loosest: an active-set method returns an exact KKT point, a splitting
method stops at a tolerance, and an interior-point method approaches the boundary
without reaching it. Every solver is therefore asked for $\qpCmpTol$ and the
residual it \emph{achieves} is reported beside its time, in the scaled sup-norm
of~\eqref{eq:certificate}. Multipliers are recovered from the returned iterate by
least squares and not read from the solver, so none is charged for reporting
them in another convention.

\begin{table}[ht]
\centering
\small
\caption{Time per solve in milliseconds and achieved KKT residual, at
$n \in \{\qpCmpSizes\}$. The first four rows are active-set methods and the last
two are not, which the residual column is there to show.}
\label{tab:compare}
\begin{tabular}{lrrrrrrrr}
\toprule
& \multicolumn{2}{c}{$n = 50$} & \multicolumn{2}{c}{$n = 100$}
& \multicolumn{2}{c}{$n = 200$} & \multicolumn{2}{c}{$n = 400$} \\
\cmidrule(lr){2-3}\cmidrule(lr){4-5}\cmidrule(lr){6-7}\cmidrule(lr){8-9}
Solver & ms & resid. & ms & resid. & ms & resid. & ms & resid. \\
\midrule
\quadprogCompareRows
\bottomrule
\end{tabular}
\end{table}

\begin{figure}[ht]
\centering
\includegraphics[width=0.86\textwidth]{graphs/quadprog_compare.pdf}
\caption{Time per solve against $n$ on three constraint shapes. The reference
implementation of the same algorithm wins at the small end, where cost is
interpreter dispatch and not arithmetic, and loses at the large end, where it
is arithmetic.}
\label{fig:compare}
\end{figure}

\paragraph{The same algorithm in two languages crosses over.} The C reference is
the fastest solver here at $n = 50$ and among the slowest at
$n = \qpCmpNhi{}$, where it is $\qpCmpVsRef\times$ slower than the
implementation measured in this paper. Nothing algorithmic separates them: they
walk the same active sets in the same order. What separates them is that below a
few hundred variables a solve costs what its operations cost to dispatch, and
above it a solve costs arithmetic, which reaches vendor-tuned kernels in one and
hand-written scalar loops in the other. This reproduces, independently and as a
by-product, the crossover the implementation is designed around.

\paragraph{Accuracy is where the interior-point method pays.} Clarabel is asked
for $\qpCmpTol$ and delivers between $10^{-9}$ and
$\qpCmpResidClarabel$, against $\qpCmpResidExact$ for the active-set method,
several orders of magnitude, and not a defect of the solver. The optima here are
vertices of the feasible polyhedron, with many constraints active, and a vertex is
exactly the case an interior-point method approaches asymptotically. OSQP reaches
$\qpCmpResidOsqp$, comparable to the active-set method, but only with its
polishing step enabled and its tolerance set to $\qpCmpTol$; at its default
tolerance it is faster and not comparable. The general point is that on
combinatorially structured optima an active-set method returns an answer the
others can only approach, and comparing times without residuals hides that.

\paragraph{A different factorisation of the same walk crosses over too.} Both
in-family solvers behave the same way. Both are fastest of everything at $n = 50$,
the C reference on bounds and on budget-plus-bounds and DAQP on dense $C$, and both
are behind at
$n = \qpCmpNhi{}$, where DAQP is $\qpCmpVsDaqp\times$ slower on bounds and
$\qpCmpVsDaqpDense\times$ on dense $C$. Accuracy is indistinguishable from ours
throughout ($\qpCmpResidDaqp$ at worst), as one expects of methods in the same
family: all three terminate on a vertex rather than approaching one, which is
exactly what separates them from the two rows below. The comparison is not
apples to apples in one respect: DAQP is written in library-free C
for embedded targets and pays no interpreter cost, so at small $n$ it is measuring
a different thing from us. What it shares is the shape of the crossover, and what
it shows is that the shape is a property of the regime and not of any one
implementation.

\paragraph{Guessing the active set is worth more than the language.} The
\texttt{fast} path is the quickest of everything measured at
$n = \qpCmpNhi{}$ on every family, and it beats the exact walk by roughly a
factor of two to three, more than the gap between the two implementations of the
exact walk at the same size. Iterations, not the cost of an iteration, dominate
at these sizes; Section~\ref{sec:pdas}'s argument arrives here from a different
direction. This ordering is the well-conditioned one, and it reverses:
Section~\ref{ssec:transform} locates the conditioning at which the certificate
begins rejecting the guess, past which the \texttt{fast} path is the slowest of
the four in-family rows rather than the quickest.

\paragraph{What this comparison is not.} These are dense problems of small to
moderate size whose optima are vertices, which is the territory
Section~\ref{sec:intro} claims for the method and precisely the territory OSQP and
Clarabel are not built for. Both are designed for large sparse problems, both
exploit sparsity this experiment does not give them, and both would be expected to
win outright at $n$ in the tens of thousands with a sparse $G$. The table is
evidence that the method is the right choice on its own ground, and no evidence at
all about anyone else's.

\paragraph{So which of the four should a reader use?} The question is fair, and the
table answers it only if one reads the columns against each other, so we state it.
For dense problems whose optima are vertices, any of the four in-family rows returns
the same answer to machine precision, and the choice is about constraints other than
accuracy. Below a few hundred variables the compiled implementations win, and by a
margin no amount of array-language tuning will close, because what is being measured
there is interpreter dispatch. Above that the ordering reverses, and the fast path of
Section~\ref{sec:pdas} leads on every family measured. The two solvers outside the
family are the right choice once $n$ is large and $G$ is sparse, which this
experiment does not test and does not claim to.

DAQP deserves a sharper statement than the ordering gives it, because on three
counts it is simply ahead: it handles semidefinite $G$ through proximal
regularisation, its worst-case iteration count can be certified offline for a
parametric family~\cite{arnstrom2022certification}, and it is library-free C sized
for embedded targets. A reader with a toolchain, a real-time deadline and a
semidefinite Hessian should use it. What the implementation studied here offers
against that is narrower: no compiler and no build step, a
drop-in signature for code already written against the reference, and the arithmetic
advantage above a few hundred variables. Whether that is the better trade is a
property of the reader's situation and not of the method.

\subsection{What these experiments do not establish}
\label{ssec:limits}

These experiments verify agreement between the code and the identities proved
here, not the absence of defects elsewhere, in constraint selection or the
infeasibility verdict or degenerate input, for which a differential comparison
of complete trajectories against an independent implementation would be the
instrument. Conditioning is no longer among the gaps, since
Section~\ref{ssec:transform} varies it over ten orders of magnitude at a fixed
answer, but it is varied there and nowhere else: every other measurement here
draws $G = BB\T + nI$ and so reports one conditioning. That choice also
biases Section~\ref{ssec:pdas}, in the direction that works against the conclusion
drawn there, since a well-conditioned $G$ spreads the minimiser so few bounds bind
where a realistic covariance concentrates it and binds many;
Section~\ref{sec:portfolio} measures how far off that is. And the checks are dense
and of modest size, with single-machine timings.
